\centerline{\bf \Huge Перестановки}


\textbf{Определение.} Рассмотрим натуральные числа $1,2, \ldots, n$. \textit{Перестановкой} чисел $1,2, \ldots, n$ называется набор тех же чисел $x_1, x_2, \ldots, x_n$, записанных в каком-то порядке. Перестановки обозначаются греческими буквами из середины алфавита: $\sigma, \pi, \tau$. Множество всех перестановок длины $n$ называется \textit{группой перестановок} (или же \textit{симметрической группой}) и обозначается через $S_n$. Ясно, что $\left|S_n\right|=n!$. Чтобы записать перестановку, применяется один из следующих способов.

\textbf{1.} \textit{Табличка}. $\sigma=\left(\begin{array}{cccc}1 & 2 & \ldots & n \\ x_1 & x_2 & \ldots & x_n\end{array}\right)$.

\textbf{2.} \textit{Ориентированный граф}. Возьмем на плоскости $n$ точек, занумеруем их числами от 1 до $n$ и направим ребра $k \rightarrow x_k$ для всех $k$ от 1 до $n$. Заметим, что данный граф - это просто объединение циклов: $\sigma=C_1 \cup \ldots \cup C_h$ (возможно, цикл имеет длину 1 , т.е. является петлей).

\textbf{3.} \textit{Диаграмма Юнга}. Сопоставим перестановке $\sigma$ диаграмму Юнга, взяв в качестве длин строк длины циклов $C_1, \ldots, C_h$.

Пусть $x_1, x_2, \ldots, x_n$ - перестановка чисел $1,2, \ldots, n$. Инверсией называется пара чисел ( $x_i, x_j$ ), такая, что $i<j$, но $x_i>x_j$. Перестановка $x_1, x_2, \ldots, x_n$ называется \textit{чётной}, если в ней чётное число инверсий, и \textit{нечётной} в противном случае. Знак перестановки $\sigma$ - это величина $\operatorname{sgn}(\sigma)$, равная +1 , если перестановка $\sigma$ четна, и -1 , если нечетна. Транспозицией называется перестановка, в которой два элемента меняются местами, а остальные - неподвижны.


\problem{1}
Рассмотрим перестановку $\sigma=\left(\begin{array}{cccccccccc}1 & 2 & 3 & 4 & 5 & 6 & 7 & 8 & 9 & 10 \\ 9 & 5 & 3 & 6 & 4 & 2 & 7 & 10 & 1 & 
8\end{array}\right)$. Изобразите соответствующий ориентированный граф и диаграмму Юнга, а также вычислите знак перестановки $\sigma$.

\problem{2}
Сколько существует различных перестановок, имеющих диаграмму Юнга, как у перестановки из предыдущего номера?

\problem{3}
Композиция перестановок. Перестановку $\sigma=\left(\begin{array}{cccc}1 & 2 & \ldots & n \\ x_1 & x_2 & \ldots & x_n\end{array}\right)$ можномыслить как функцию, определенную на множестве $\{1,2, \ldots, n\}$, полагая $\sigma(k)=x_k$ для всех $k=1$, ..., n. В таком случае композицией перестановок называется их последовательное применение: $(\tau \circ \sigma)(k)=\tau(\sigma(k))$.

(a) Вычислите композицию перестановок $\sigma=\left(\begin{array}{cccccccccc}1 & 2 & 3 & 4 & 5 & 6 & 7 & 8 & 9 & 10 \\ 9 & 5 & 3 & 6 & 4 & 2 & 7 & 10 & 1 & 8\end{array}\right)$ и $\tau=\left(\begin{array}{lllllllccc}1 & 2 & 3 & 4 & 5 & 6 & 7 & 8 & 9 & 10 \\ 3 & 8 & 7 & 1 & 5 & 2 & 4 & 10 & 6 & 9\end{array}\right)$.

(b) Докажите, что для каждой перестановки $\sigma$ существует такая перестановка $\tau$, что $(\tau \circ \sigma)(k)=k$ для всех $k=1, \ldots, n$. Такая перестановка называется обратной $\kappa \sigma$ и обозначается через $\sigma^{-1}$.

(c) Известно, что у перестановки $\sigma$ все циклы имеют нечетную длину. Докажите, что существует такая перестановка $\rho$, что $\rho \circ \rho=\sigma$.

(d*) Пусть $\sigma$ и $\tau$ - произвольные перестановки. Докажите, что у перестановок $\sigma$ и $\tau^{-1} \circ \sigma \circ \tau$ одинаковые диаграммы Юнга.

\problem{4}

(a) Чему равен знак перестановки $(n, n-1, n-2, \ldots, 2,1)$ ?

(b) Докажите, что транспозиция двух элементов перестановки меняет ее четность.

(c) Дана перестановка чисел $1, \ldots, n$, где $n$ число вида $4 k+2$. Разрешается менять подряд идущие числа $i, j, \ell$ на числа $j, \ell, i$. Можно ли такими операциями получить из перестановки $1,2, \ldots, n$ перестановку $n, n-1, \ldots, 1$ ?

(d) Докажите, что четных и нечетных перестановок поровну.

\problem{5}
На физкультуре $n$ школьников выстроились в порядке обратном их росту. Придя на урок, учитель попросил выстроиться ученикам, как и принято, то есть по росту. Во избежании суеты, разрешается поменяться двум соседним ученикам. Сколько обменов необходимо произвести школьникам, чтобы добиться требуемого?

\problem{6}
\textbf{Задача Константинова об обмене квартир.}
В некотором городе разрешаются только парные обмены квартир (если две семьи обмениваются квартирами, то в тот же день они не имеют права участвовать в другом обмене). Докажите, что любой обмен квартирами можно осуществить за два дня.